\documentclass{article}

\textwidth=6.5in
\textheight=8.5in
\voffset=-54pt
\oddsidemargin=0in
\evensidemargin=0in
\setlength{\parskip}{8pt}
\setlength{\parindent}{0pt}

\usepackage{workbook}

\usepackage[sols]{handout}

\begin{document}

\begin{center}
  \large\textsc{Problem Set 8 - Sinusoidal Functions}
  
      \pspace{12pt}
  \begin{problemonly}\large Name: \rule{5cm}{0.15mm}\\ \end{problemonly}
\end{center}

\begin{grading}
\begin{enumerate}
    \item 10
    \item 5
    \item 10
    \item 15
    \item 12
\end{enumerate}
Total: 52 points.
\end{grading}

 Pages 603 - 609 in Gottlieb can help you with the
  following problems.

\begin{enumerate}[topsep=0pt]
\item 

\begin{grading}
10 points. 6 for part (a), 1 for part (b), 2 for part (c), and 1 for part (d). 
\end{grading}

%({Problem Set 8}, {{\#1}})
  \begin{enumerate}
  \item
\begin{problem}
      On the same set of axes, graph the functions $y=\cos (2t)$ and $y=2 \cos
      t$ on the interval $[0, 2 \pi]$.  
    \end{problem}

        \begin{tikzpicture}
          \begin{axis}[xlabel=$t$, ylabel=$y$,
          xmin=-0.5, xmax=6.5,
          xtick={0, 0.785, ..., 6.281}, 
          xticklabels={0, , $\pi/2$ , ,$\pi$, , $3\pi/2$, ,$2 \pi$}, 
          ytick={-2, -1, ..., 2}, 
          width=6in,
          height=3.5in
          ]
            \addplot+[domain=0:13, draw=none]{cos(deg(2*x))};
            \addplot+[domain=0:13, draw=none]{2*cos(deg(x))};
            \solonly{
            \addplot+[domain=0:13, red]{cos(deg(2*x))};
            \addplot+[domain=0:13, blue]{2*cos(deg(x))};
            }
          \end{axis}
        \end{tikzpicture}

    \begin{solution}
      $y=\cos(2t)$ is the red graph, a horizontal scaling of $y = \cos t$ by a
      factor of $\frac{1}{2}$, and $y=2 \cos t$ is the blue graph, a vertical
      scaling of $y = \cos t$ by a factor of $2$.
    \end{solution}

  \item
    \begin{problem}[\stretch{0.5}]
      True or false: $\cos(2x) = 2 \cos x$ for all $x$.
    \end{problem}

    \begin{solution}
      False.  (It's clear that the graphs we drew in
      {{{{\#1}}(a)}} are different.)
    \end{solution}

  \item

    \note{This is here mostly because I want them to get used to the idea
      of looking at graphs to determine how many solutions an equation
      has.}

    \begin{problem}[\stretch{1}]
      Based on your graphs, how many solutions does
      $\cos(2x) = 2 \cos x$ have on the interval $[0, 2 \pi]$?
    \end{problem}

    \begin{solution}
      The graphs intersect at two places, so the equation should have
      \fbox{two} solutions.
    \end{solution}

  \item
    \begin{problem}[\nstretch{1}]
      How many solutions does the equation $\cos(2x) = 2 \cos x$ have on
      the interval $\left[0, \frac{\pi}{2} \right]$?
    \end{problem}

    \begin{solution}
      \fbox{Zero}. The graphs do not intersect on that interval.
    \end{solution}

  \end{enumerate}

\item 

\begin{grading}
5 points. Be sure that they have plotted the graph on the correct domain. Also, when labeling a local min and max, they must label the points with the exact coordinates! They cannot simply highlight where they occur on the graph. 
\end{grading}

%({Problem Set 8}, {{\#2}})

\begin{problem}
In this problem, we are going to sketch the graph of $y=-2 \sin ( \frac{x}{3})
    + 1$.
\end{problem}

\begin{enumerate}

\item
\begin{problem}
    Find the amplitude, midline, and period.
\end{problem}
\pspace{\stretch{1}}

\begin{solution}
    \begin{itemize}
        \item Amplitude $=2$
        \item Midline: $y=1$
        \item Period $=6\pi$
    \end{itemize}
\end{solution}

\item
\begin{problem}[\stretch{3}]
    Sketch the graph on the interval $[0, 12 \pi]$.  Label the coordinates of a local
    minimum, a local maximum, and the $y$-intercept.
  \end{problem}

  \begin{solution}
    We can get this graph by starting with $y = \sin x$, flipping
    vertically and stretching vertically by a factor of $2$, stretching
    horizontally by a factor of $3$ (so that the period is $\frac{2
      \pi}{1/3} = 6 \pi$), and shifting up by 1:
    \begin{center}
      \begin{tikzpicture}
        \begin{axis}[xtick={0., 9.42478, 18.8496, 28.2743,
            37.6991}, xticklabels={$0.$, $3 \pi$, $6 \pi$, $9 \pi$, $12
            \pi$}, ytick={-1, 1, 3}, clip=false, width=3in]
          \addplot+[domain=0:38]{-2*sin(deg(x)/3)+1};
          \filldraw (axis cs: 0, 1) circle(2pt) node[anchor=south west,
          inner sep=1pt]{$(0, 1)$};
          \addplot[closed, mark size=2] coordinates {(0,1) (3*pi/2, -1) (9*pi/2, 3) (15*pi/2, -1) (21*pi/2,3)};
          \node[below] at (3*pi/2, -1) {$(\frac{3\pi}{2}, -1)$};
          \node[above] at (9*pi/2, 3) {$(\frac{9\pi}{2}, 1)$};
          \node[below] at (15*pi/2, -1) {$(\frac{3\pi}{2}, -1)$};
          \node[above] at (21*pi/2, 3) {$(\frac{3\pi}{2}, 1)$};
        \end{axis}
      \end{tikzpicture}
    \end{center}
  \end{solution}

\end{enumerate}
\item 

\begin{grading}
10 points. 5 for each graph. Students may have chosen to use sine for the first graph, or cosine for the second graph. If their equations still match the graph, they should get full credit, but you should highlight why their choice wasn't the most efficient. (See \#5d.)
\end{grading}
%({Problem Set 8}, {{\#3}})

Find equations that fit the curves below.


  \begin{multicols}{2}
    \begin{enumerate}

    \item
      \raisebox{2ex-\height}{
        \begin{tikzpicture}
          \begin{axis}[xtick=\empty, ytick=\empty, clip=false, width=2in,
            height=1.5in, enlarge y limits=0.2] 
            \addplot+[domain=-15:15, semithick]{2*pi*cos(deg(x)/2)+pi};
            \pgfmathsetmacro\x{3*pi};
            \pgfmathsetmacro\y{2*pi};
            \pgfmathsetmacro\z{-pi};
            \filldraw (axis cs: 0, \x) circle(1.5pt) node[anchor=south
            west]{$(0, 3\pi)$};
            \filldraw (axis cs: \y, \z) circle(1.5pt) node[anchor=north
            west]{$(2\pi, -\pi)$};
          \end{axis}
        \end{tikzpicture}
      }

      \begin{solution}
        This graph has amplitude $2 \pi$, balance value $\pi$, and period
        $4 \pi$.  It has a local maximum at $x = 0$, so we can get this
        graph by starting with $y = \cos x$, stretching vertically by a
        factor of $2 \pi$, stretching horizontally by a factor of $2$, and
        shifting up by $\pi$.  This gives us the equation 
        $\boxed{y = 2 \pi \cos \left( \frac{x}{2} \right) + \pi}$.

          \checkmark Let's check by plugging in $x = 0$ and $x = 2 \pi$ (the
          $x$-coordinates of the two points we're given): when $x = 0$, $y
          = 2 \pi \cos(0) + \pi = 3 \pi$; when $x = 2 \pi$, $y = 2 \pi
          \cos (\pi) + \pi = - \pi$, which agrees with the two points we
          were given. 

      \end{solution}

      \columnbreak

    \item
      \raisebox{2ex-\height}{
        \begin{tikzpicture}
          \begin{axis}[xtick=\empty, ytick=\empty, clip=false, width=2in,
            height=1.5in, enlarge y limits=0.2] 
            \addplot+[domain=-3:3, semithick]{-2*sin(pi*deg(x))+3};
            \filldraw (axis cs: -2.5, 5) circle(1.5pt) node[anchor=south]{$\left( - \frac{5}{2}, 5 \right)$};
            \filldraw (axis cs: 0.5, 1) circle(1.5pt) node[anchor=north,
            fill=white, yshift=-2pt]{$\left( \frac{1}{2}, 1 \right)$};
          \end{axis}
        \end{tikzpicture}
      }

      \begin{solution}
        This graph has amplitude $2$, balance value $3$, and period $2$
        (there are three half-periods between $x = - \frac{5}{2}$ and $x =
        \frac{1}{2}$).  At $x = 0$, the graph is at its balance value and
        decreasing, so it seems easiest to model this by starting with $y
        = - \sin x$.  If we stretch this vertically by a factor of 2,
        stretch it horizontally by a factor of $\frac{1}{\pi}$, and shift
        it up by $3$, we'll end up with the graph we want.  Therefore, one
        model is $\boxed{y = -2 \sin (\pi x) + 3}$.

        \checkmark
          Let's check by plugging in $x = - \frac{5}{2}$ and $x =
          \frac{1}{2}$.  When $x = - \frac{5}{2}$, $y = - 2 \sin \left( -
            \frac{5 \pi}{2} \right) + 3 = -2(-1) + 3 = 5$; when $x =
          \frac{1}{2}$, $y = -2 \sin \left( \frac{\pi}{2} \right) + 3 =
          1$, which agrees with the points we're given.
        
      \end{solution}

    \end{enumerate} 
    \end{multicols}

  \probonly{\emph{Check your answers by using a calculator or computer to
      graph your equations.  }}

  \pspace{\nstretch{2}}

\item 

\begin{grading}
15 points. 5 for part (a), 5 for part (b), 1 for part (c), 1 for part (d)i, and 3 for part (d)ii.
\end{grading}

%({Problem Set 8}, {{\#4}})

Between 2008 and 2014, the world's largest Ferris wheel was the
  Singapore Flyer. It has a diameter of 150 meters and makes one
  counterclockwise revolution every $30$ minutes.  At time $t = 0$, you
  are in a car at the bottom of the wheel, which is a height of 15 meters
  above the ground.

  \begin{enumerate}[topsep=0pt]
  \item
    \begin{problem}[\stretch{2}]
      Let $h(t)$ be your height (in meters) above the ground at time $t$,
      where $t$ is measured in minutes.  Graph $h(t)$ from time $t=0$ to
      time $t=90$.
    \end{problem}

    \begin{solution}
      At time 0, you are at a height of 15.  Your height then increases
      until you reach the top of the Ferris wheel, which has a height of
      165 (your initial height of 15 meters plus the diameter of the
      Ferris wheel).  This takes 15 minutes since it makes up half a turn
      of the Ferris wheel.  You then come back down until, at time 30,
      your height is again 15 meters.  This repeats over and over, giving
      a graph like this:
      \begin{center}
      \begin{tikzpicture}
        \begin{axis}[xtick={30, 60, 90}, ytick={15, 90, 165}, xlabel=$t$,
          ylabel=$h(t)$, width=6in, height=2in,
              enlarge x limits=0.05]
          \addplot+[domain=0:90] {-75*cos(2*pi/30*deg(x))+90};
          \addplot[dotted, domain=0:90] {90};
        \end{axis}
      \end{tikzpicture}
    \end{center}
    \end{solution}

  \item
    \begin{problem}[\stretch{1}]
      Write an explicit formula for $h(t)$.  (You may want to use a calculator or
      computer to check that your equation matches your graph.)
    \end{problem}

    \begin{solution}
      The graph has amplitude $75$, midline value $90$, and period $30$.
      Based on our sketch, it's probably easiest to model with the
    negative of a cosine function, since the function we're modeling has
    a minimum at $t = 0$.  We can get our function by starting with $-
    \cos t$, stretching vertically by a factor of 75 (to get the right
    amplitude), stretching horizontally by a factor of $\frac{30}{2
      \pi}$ (to get the right period), and shifting up by 90 (to get the
    right midline value).  So, our model is $h(t) = -75 \cos \left(
      \frac{2 \pi}{30} t \right) + 90$, which simplifies to $\boxed{h(t)
      = -75 \cos \left( \frac{\pi}{15} t \right) + 90}$.

      \checkmark
    It's easy to get confused about horizontal stretches, so it's a very
    good idea to plug in a few points to check your answer.  Here, let's
    plug in $t = 0$ and $t = 15$ to see if they give us what we expect:
    \begin{align*}
      h(0) &= -75 \cos(0) + 90 = -75 + 90 = 15 \text{ \checkmark} \\
      h(15) &= -75 \cos (\pi) + 90 = 75 + 90 = 165 \text{ \checkmark}
    \end{align*}
  
    \end{solution}

  \item
    \begin{problem}[\stretch{0.5}]
      At time $t = 50$, is your height above the ground increasing or
      decreasing? 
    \end{problem}

    \begin{solution}
      From our graph, we see that it is \fbox{decreasing}.  (It increases
      from $t = 30$ to $t = 45$, then decreases from $t = 45$ to $t =
      60$.) 
    \end{solution}

  \item
    \begin{enumerate}
    \item 
      \begin{problem}[\stretch{0.5}]
        At how many different times between $t = 0$ and $t = 90$ is your
        height above the ground exactly 50 meters? 
      \end{problem}

      \begin{solution}
        From our graph, we can see that there are \fbox{6} such times. 
      \end{solution}

    \item 
      \begin{problem}[\stretch{1.5}]        
        If the \textbf{first} of these times is called time $t_0$, express
        the other times in terms of $t_0$.  Label these on your graph. 
      \end{problem}

      \begin{solution}
        \begin{center}
          \begin{tikzpicture}
            \begin{axis}[xtick={4.81409, 25.1859, 64.8141, 55.1859,
                85.1859, 34.8141, 30, 60, 90}, xticklabels={ $t_0$, $30-t_0$,
                $60+t_0$, $60-t_0$, $90-t_0$, $30+t_0$}, ytick={15, 50, 90,
                165}, xlabel=$t$, ylabel=$h$, width=6in, height=2in,
              enlarge x limits=0.05]              
              \addplot+[domain=0:90] {-75*cos(2*pi/30*deg(x))+90};
              \addplot+[domain=0:90, dashed, black] {50};
              \addplot[domain=0:90, dotted] {90};
              \addplot[closed] coordinates {(4.81, 50) (25.19, 50) (64.81, 50) (55.19,50) (85.19, 50) (34.81, 50)};
            \end{axis}
          \end{tikzpicture}
        \end{center}
        Using the graph, we can see that the second time the height is 50
        is $t_0$ minutes before $t = 30$, or $t = 30 - t_0$.  The next
        pair of times that the height is 50 happen exactly 30 minutes (one
        period) after the first pair, so they are $30 + t_0$ and $30 + (30
        - t_0) = 60 - t_0$.  The next pair of times happens after another
        30 minutes, so they are $30 + (30 + t_0) = 60 + t_0$ and $30 + (60
        - t_0) = 90 - t_0$.  So, the 5 times after $t_0$ when the height
        is 50 are \fbox{$30 - t_0$, $30 + t_0$, $60 - t_0$, $60 + t_0$,
          and $90 - t_0$}.
      \end{solution}

    \end{enumerate}
  \end{enumerate}


\probonly{\newpage}



\item 

\begin{grading}
12 points. 4 for part (a), 4 for part (b), 2 for part (c), and 2 for part (d).
\end{grading}

\textbf{Reflect Back.} \#1(b) shows us that, when manipulating trigonometric expressions, we need to carefully consider which rules aka identities are actually valid. In each of the following parts, explain why the equation is a valid identity (true for all possible values), or produce a counterexample.

You should \textit{not} need a calculator for these problems!

\begin{enumerate}
    \item 
    \begin{problem}[\vfill]
    Show that \[\cos(A+B) \stackrel{?}{=} \cos(A)+\cos(B)\] is not a valid identity by finding specific values for $A$ and $B$ that make the equation false. (In other words, find a \textit{counterexample} to the statement.)

    Do the same for $\sin(A+B) \stackrel{?}{=} \sin(A)+\sin(B)$.
    \end{problem}

    \begin{solution}
        To see that $\cos(A+B)\ne \cos(A)+\cos(B)$ in general, we can let $A=B=0$. Then:
        \begin{align*}
            \cos(A+B) &= \cos\left(0+0\right)=\cos(0)=\boxed{1}\\
            \cos(A)+\cos(B) &= \cos\left(0\right)+\cos\left(0\right)=1+1=\boxed{2}
        \end{align*}
        
        To see that $\sin(A+B)\ne \sin(A)+\sin(B)$ in general, we can let $A=B=\frac{\pi}{2}$. Then:
        \begin{align*}
            \sin(A+B) &= \sin\left(\frac{\pi}{2}+\frac{\pi}{2}\right)=\sin(\pi)=\boxed{0}\\
            \sin(A)+\sin(B) &= \sin\left(\frac{\pi}{2}\right)+\sin\left(\frac{\pi}{2}\right)=1+1=\boxed{2}
        \end{align*}

        
    \end{solution}

    \item 
    \begin{problem}[\vfill]
    Show that \[\cos(AB)\stackrel{?}{=} \cos(A)\cos(B)\] is not a valid identity by finding a counterexample. (Find specific values for $A$ and $B$ that make the equation false.)

    Do the same for the equation $\sin(A\cdot B)\stackrel{?}{=} \sin(A)\cdot\sin(B)$.
    \end{problem}

    \begin{solution}
        To see that $\cos(A B)\ne \cos(A)\cos(B)$ in general, we can let $A=0$ and $B=\pi$. Then:
        \begin{align*}
            \cos(AB) &= \cos(0\cdot\pi)=\cos(0)=\boxed{1}\\
            \cos(A)\cos(B) &= \cos(0)\cos(\pi)=(1)(-1)=\boxed{-1}
        \end{align*}
        
        To see that $\sin(A B)\ne \sin(A)\sin(B)$ in general, we can let $A=B=\pi$. Then:
        \begin{align*}
            \sin(A B) &= \sin(\pi\cdot\pi)=\sin({\pi}^2)\ne 0\\
            \sin(A)\sin(B) &= \sin(\pi)\sin(\pi)=0\cdot 0=0
        \end{align*}
        
        \textcolor{magenta}{How do we know that $\sin({\pi}^2)\ne 0$? Well, from the unit circle, we know that $\sin(t)$ is only equal to 0 if $t$ is an integer multiple of $\pi$ (e.g., $...,-3\pi, -2\pi, -\pi, 0, \pi, 2\pi, 3\pi, ...$). However, ${\pi}^2$ is \textit{not} an integer multiple of $\pi$. So, although we don't know the value of $\sin({\pi}^2)$ exactly, we know that it can't equal 0.}

    \end{solution}

    \item 
    \begin{problem}[\vfill\newpage]
    Use graph transformations to explain why \[\cos(t)=\sin\left(t+\frac{\pi}{2}\right)\] is a valid identity (true for all values of $t$.)
    \end{problem}

    \begin{solution}
        The graph of $\sin\left(t+\frac{\pi}{2}\right)$ is just the graph of $\sin(t)$ shifted to the left by $\frac{\pi}{2}$ units. However, if we shift the graph of $\sin(t)$ to the left by exactly $\frac{\pi}{2}$ units, we can see that the graph we get is exactly the same as the graph of $\cos(t)$.
    \end{solution}

\end{enumerate}

    \item 
    \begin{problem}
    For the same reasoning as in \#5(c), any sinusoidal function written in terms of cosine can be rewritten in terms of sine (and vice-versa). Thus, when writing an equation for a sinusoidal curve, as you did in \#3, we \textit{always} have the choice to use either sine OR cosine.

    Looking back at \#3, explain in words why it is simpler to write an equation for the graph in \#3(a) using cosine, while the graph in \#3(b) is more easily expressed using sine. 
    \end{problem}
    \pspace{\stretch{2}}

    \begin{solution}
        The graph in \#3(a) has a local maximum at $x=0$, as does the graph of $\cos(t)$. So, by modeling this graph with cosine, we can avoid the need for horizontal shifts, and simply focus on getting the amplitude, period, and vertical shift correct.

        Similarly, the graph in \#3(b) has a balance point at $x=0$, as does the graph of $\sin(t)$. So, once again, we can avoid having to do any horizontal shifts if we modeled the graph using a sine function.
    \end{solution}


\item 
\begin{problem}
    There are many different trigonometric identities that can occasionally come in handy. We will not ask you to know \emph{all} of them by heart, but it's useful to know that they exist! Here are two that we may occasionally use, but you do not need to memorize them. They are sometimes called the the \emph{sum identities} or the \emph{addition formulas}.

    \begin{itemize}
        \item $\sin(A+B)=\sin(A)\cos(B)+\cos(A)\sin(B)$
        \item $\cos(A+B)=\cos(A)\cos(B)-\sin(A)\sin(B)$
    \end{itemize}

    Use the sine addition formula to show that $\sin(2t)=2\sin(t)\cos(t)$, which is one of the so-called \emph{double-angle formulas}. The first and last steps have been provided. Fill in the intermediate steps.
\end{problem}
\begin{flalign*}
    \sin(2t) &=  \solonly{\sin(t+t)} &\\[24pt]
    &= \solonly{\sin(t)\cos(t) + \cos(t)\sin(t)}&\\[24pt]
    &= 2\sin(t)\cos(t)&
\end{flalign*}



\end{enumerate}

\psreflection{Optional reading for next class is \S 19.3, 20.1 in Gottlieb.}

\end{document}